\subsection{All finite changes of an integer exponent}

The preceding erasure argument admits a substantial strengthening on
CM elliptic curves. We prove the Kummer and independence statements
needed for it, including at coefficient primes two and three. The
conclusion concerns the geometric abelian Selmer quotient, so it
also applies to the intrinsic metabelian Selmer quotient in A.12.

\begin{statement}{A.61}{CM Kummer openness and simultaneous zero-primary Frobenius}{Proved}
Let $E/L$ be a CM elliptic curve over a number field containing all
its endomorphisms, let $R\in E(L)$ have infinite order, and let
$\mathcal P$ be a nonempty finite set of rational primes. Choose
compatible division points of $R$ for each $p\in\mathcal P$, and
write
\[
 H_p=\operatorname{im}\bigl((c_p,\rho_p):
 G_L\longrightarrow T_pE\rtimes G_p\bigr),\qquad
 G_p=\rho_p(G_L).
 \tag{A61.1}
\]
Then $H_p$ is open in $T_pE\rtimes G_p$. The joint image
$H_{\mathcal P}$ is open in the product of these groups.
Both assertions remain true after any finite extension of $L$.

For each such extension $F/L$ there is a nonempty open,
conjugacy-invariant subset of the joint image, determined in a
finite quotient, such that every good prime $w\nmid\prod_{p\in
\mathcal P}p$ with Frobenius in that subset satisfies
\[
 e_{\mathcal P}R=0\quad\hbox{in }E(F_w).
 \tag{A61.2}
\]
Here $e_{\mathcal P}$ is the idempotent of A.60.
Every prime divisor of the order of a finite quotient of
$H_{\mathcal P}$ belongs to
\[
 C_{\mathcal P}
  =\bigcup_{p\in\mathcal P}
       \{\lambda:\lambda\mid p(p^2-1)\}.
 \tag{A61.3}
\]
The same bound holds for closed subgroups, in particular for the
image after finite base extension.
\end{statement}
\begin{proof}
Put $K=\operatorname{End}_L(E)\otimes\Q$, an imaginary quadratic
field, and $V=T_pE\otimes\Q_p$. The algebra $K_p=K\otimes\Q_p$
acts faithfully on $V$, which is a free $K_p$-module of rank one.
Thus $G_p$ is contained in the commutative group $K_p^\times$.

Let $J_p\subset T_pE$ be the image of $\ker\rho_p$ under $c_p$.
It is a closed additive $\Z_p$-submodule, stable under $G_p$.
Write $W=\Q_p J_p\subset V$. Choose a good rational prime
$q\ne p$ splitting completely in $L$, and choose a place of $L$
above it. Its Frobenius $A\in G_p$ has integral trace and
determinant $q$. It is not scalar: if $A=tI$, its trace would give
$t\in\Q$ and its determinant would give $t^2=q$, impossible.
Consequently $\Q_p[A]=K_p$, including when $p$ splits or ramifies
in $K$. It follows that $W$ is a $K_p$-submodule. Also
\[
 \det(A-I)=\#E(\mathbf F_q)\ne0.
 \tag{A61.4}
\]
These observations use a degree-one Frobenius, not the assumption
that $E[p]$ is nonrational.

Suppose $W\ne V$. Since $K_p$ is a product of fields, a nonzero
idempotent $e\in K_p$ selects a quotient component with
$eW=0$ and $eV\ne0$; in the nonsplit case take $e=1$.
The cocycle $ec_p$ vanishes on $\ker\rho_p$ and hence descends
continuously to $G_p$. For a cocycle $b$ of a commutative group
$G$ and a central element $A\in G$, the cocycle identity gives
\[
 (A-1)b(g)=(g-1)b(A).
\]
If $A-1$ is invertible on the coefficient space, this says
$b(g)=(g-1)(A-1)^{-1}b(A)$, a coboundary. Apply this identity
on $eV$ using (A61.4). We obtain $e\kappa_p(R)=0$ in
$H^1(L,eV)$.

This contradicts the rational Tate Kummer injection
\[
 E(L)\otimes\Q_p\lhook\joinrel\longrightarrow H^1(L,V).
\]
Indeed the map $K\to E(L)\otimes\Q$, $\alpha\mapsto\alpha R$,
is injective: a nonzero CM endomorphism is an isogeny and cannot
send $R$ to torsion. Tensoring with $\Q_p$ preserves that
injection, so $eR\ne0$ for every nonzero CM idempotent $e$.
The injection used here follows directly from the multiplication
Kummer sequences and Mordell--Weil; passage from $T_pE$ to $V$
kills only torsion. Hence $W=V$, so $J_p$ has finite index in
$T_pE$. The projection $H_p\to G_p$ is surjective, and the
cosets of $H_p$ in $T_pE\rtimes G_p$ are represented by
$T_pE/J_p$. Thus $H_p$ is open. This proof applies without
change over every finite extension of $L$.

We justify the joint assertion, since separately open projections
alone would not suffice for arbitrary profinite groups. Choose
normal open pro-$p$ subgroups $N_p\subset H_p$, for instance
inside the kernels of the affine reductions modulo $p$. If $H$
is the joint image, then
$H_0=H\cap\prod_pN_p$ has finite index in $H$. Its projection
$N'_p$ to each $N_p$ is therefore open. A closed subgroup of a
finite product of pro-$p$ groups for distinct primes is the product
of its projections. To see this, for $x=(x_p)$ in the subgroup,
choose integers tending to $1$ in $\Z_p$ and to $0$ in all the
other $\Z_{p'}$. Taking these powers of $x$ isolates its $p$th
component; closedness places that component in the subgroup.
Multiplying the isolated components proves the assertion.
Thus $H_0=\prod_pN'_p$, which is open in $\prod_pH_p$.
This proves joint openness, also after any finite base extension.

Fix such an extension $F$. Joint openness provides a product of
identity neighbourhoods contained in the joint image, with an
open subgroup $G_{p,0}$ in each torsion factor. Choose
$A_p\in G_{p,0}$ with $A_p-I$ invertible. Such an element exists
even if $E[p]$ is rational over $F$: choose a good rational prime
splitting completely in the finite extension cutting out
$G_{p,0}$ and use its degree-one Frobenius, as in (A61.4).
The tuple $((0,A_p))_p$ belongs to the joint image. For each $p$
choose $b_p$ such that
$p^{b_p}(A_p-I)^{-1}$ is integral. On a sufficiently small
congruence neighbourhood of this tuple one has
\[
 A'_p\equiv A_p\pmod{p^{b_p+1}},\qquad
 c'_p\equiv0\pmod{p^{b_p}},\qquad
 c'_p\in(A'_p-I)T_pE.
 \tag{A61.5}
\]
The last implication follows by factoring
$A'_p-I=(A_p-I)(1+(A_p-I)^{-1}(A'_p-A_p))$; the second factor
is integral and invertible. The condition in (A61.5), together
with nonzero determinant, is invariant under affine conjugacy.
Taking the conjugacy saturation of this finite congruence
neighbourhood gives the required nonempty open subset.

At a good prime $w\nmid p$, the local Kummer class of $R$ is its
Frobenius cocycle in
\[
 T_pE/(A'_p-I)T_pE
   \simeq E(\mathbf F_w)[p^\infty].
\]
The formal kernel of good reduction is pro-residue-characteristic,
so this is also the $p$-primary factor of $E(F_w)$. Condition
(A61.5) makes this factor of $R$ zero for every $p\in\mathcal P$.
The division towers are unramified at such good primes: the
multiplication maps and the fibres over the integral section $R$
are finite etale over the local ring. This proves (A61.2).

Finally the kernel of affine reduction modulo $p$ is pro-$p$,
and the finite affine group has order
\[
 p^2\lvert\mathrm{GL}_2(\mathbf F_p)\rvert
    =p^3(p-1)^2(p+1).
\]
For the finite product, every closed subgroup has a normal open
subgroup whose finite quotients have prime divisors in
$\mathcal P$, and the remaining finite quotient embeds in the
product of these affine groups. This proves (A61.3), with no
omitted small-prime exception.
\end{proof}

\begin{statement}{A.62}{A finite set of altered coefficient primes forces all other unit characters}{Proved}
In the notation of A.61 let $D\subset E$ be a finite reduced
divisor, $U=E-D$, and let an abelian Selmer section of $U/L$ have
compact Kummer class $a\kappa(R)$. Suppose that $n\in\Z$ and
\[
 a_p=n\quad\hbox{for every }p\notin\mathcal P.
 \tag{A62.1}
\]
Put $B=nR$. Let $f\in L[U]^\times$ extend to a regular nonzero
function at $B$, including when $B\in D$. Write $\beta_f$ for
the difference between the global completed Kummer class obtained
from $f$ and $\kappa(f(B))$. If
\[
 \lambda\notin C:=C_{\mathcal P}\cup\{2\}
       \cup\{\hbox{rational primes ramified in }L\},
 \tag{A62.2}
\]
then the entire $\lambda$-adic component of $\beta_f$ is zero.
This assertion imposes no finite-ramification condition on the
original completed Kummer class.
\end{statement}
\begin{proof}
Fix $r\ge1$ and let $\delta$ be the reduction of $\beta_f$ in
$H^1(L,\mu_{\lambda^r})$. Put $F=L(\mu_{\lambda^r})$.
Restriction to $F$ is injective in this $H^1$. Indeed $L$ is
unramified at $\lambda$, whereas every nontrivial subextension
of $\Q(\mu_{\lambda^r})/\Q$ is totally ramified there, so
\[
 \operatorname{Gal}(F/L)=(\Z/\lambda^r)^\times.
\]
The inflation--restriction kernel is
$H^1(\operatorname{Gal}(F/L),\mu_{\lambda^r})$.
This is zero by the elementary central-element identity in
A.61, using the cyclotomic automorphism $2$: it acts on the
additive module by $2$, whose difference from $1$ is invertible.
We excluded $\lambda=2$ explicitly.

Suppose $\delta\ne0$. Over $F$ it defines a nontrivial cyclic
extension $M/F$ of $\lambda$-power degree. Let $\mathcal T/L$
be the joint division field of A.61. Every finite quotient of
$\operatorname{Gal}(F\mathcal T/F)$ has order prime to
$\lambda$, by (A61.3); hence
\[
 M\cap F\mathcal T=F.
 \tag{A62.3}
\]
Apply A.61 \emph{after this particular base extension $F$}.
It supplies a nonempty finite Frobenius condition in
$\operatorname{Gal}(F\mathcal T/F)$ for which all the
$\mathcal P$-primary components of $R$ vanish. Choose a finite
subextension realizing that condition. By (A62.3) its Galois
group and $\operatorname{Gal}(M/F)$ occur as a direct product.
Chebotarev therefore supplies good finite places $w$ of $F$,
away from every specified finite exceptional set, whose
Frobenius satisfies the zero-primary condition and is nontrivial
in $\operatorname{Gal}(M/F)$. The zero-primary neighbourhood is
chosen over $F$, so no compatibility with the possibly nontrivial
intersection $\mathcal T\cap F$ has been assumed.

At such a place, (A62.1) and (A61.2) make the forced compact
point of the local Selmer section equal to $B$. We use full
local Tate Kummer injectivity, as in A.60. If $B\in U$, the
local section is the point section of $B$. If $B\in D$, it is
a cusp section with centre $B$. In the latter case the map
defined by $f$ extends as a unit at $B$, so it kills the cusp
inertia and sends every tangent parameter to the point section
of $f(B)$ on $\mathbf G_m$. In both cases
\[
 (\beta_f)_w=0.
\]
This contradicts the nontrivial Frobenius in the unramified
$\lambda$-Kummer extension $M/F$: its defining root could not
then lie in $F_w$. Hence $\delta=0$.

This holds for every $r$. The inverse-limit identification
$H^1(L,\Z_\lambda(1))=\varprojlim_rH^1(L,\mu_{\lambda^r})$
has no extra $\varprojlim^1$ term because the groups of rational
roots of unity at finite levels are finite. Thus the entire
$\lambda$-adic component vanishes, as asserted. Notice that
the proof only uses one finite Kummer extension at a time; its
finite ramification is automatic and does not bound the support
of the original completed class.
\end{proof}

\begin{statement}{A.63}{Every finite alteration of an integer dies at the abelian Selmer level}{Proved}
Let $E/k$ be a CM elliptic curve over a number field, let $N\ge2$,
put $U=E-E[N]$, and let $R\in E(k)$ have infinite order. If a
geometrically abelian Selmer section of $U/k$ has compact class
$a\kappa(R)$, and $a_p=n$ for all but finitely many coefficient
primes $p$, for some $n\in\Z$, then
\[
 a=n\quad\hbox{in }\widehat\Z.
 \tag{A63.1}
\]
Conversely every integer $n$ occurs, by the point section of
$nR$ when $n\ne0$ and by a cusp section at $O$ when $n=0$.

Consequently, for the rank-one cubic in A.43, the intersection of
the abelian Selmer lifting set with the set of all finite
coefficient alterations of integers is exactly $\Z$. The same
intersection is $\Z$ for the intrinsic metabelian Selmer lifting
set. This includes arbitrary changes at a single coefficient
prime, not only the erasures of A.60, and includes exponents
supported at a finite nonempty set of coefficient primes.
\end{statement}
\begin{proof}
Pass to one finite extension $L/k$ containing all CM
endomorphisms and splitting $E[N]$. A putative Selmer section
restricts, and its compact class is still $a\kappa(R)$.
Let $\mathcal P$ contain all coefficient primes at which
$a_p\ne n$. If it is empty there is nothing to prove.
Write $B=nR$.

We need one nonconstant unit $f$ on $U$ which extends as a unit
at $B$. When $n\ne0$, $B$ is nontorsion and hence belongs to
$U$; take any $f_d$ from (A60.3). When $n=0$, choose distinct
nonzero $d,e\in E[N](L)$, and take $f=f_d/f_e$.
Its divisor is $N(d)-N(e)$, so it is regular and nonzero at
$O=B$. Such $d,e$ exist for every $N\ge2$. Scaling $f$ if
desired makes $f(B)=1$ in this case; no tangent parameter enters
its value.

Apply A.62 to this $f$. At a finite place $v$ of $L$ whose
residue characteristic $q$ is outside $C\cup\mathcal P$, the
local compact point is
\[
 Q_v=aR=B+T_v,\qquad
 T_v=\sum_{p\in\mathcal P}(a_p-n)e_pR.
 \tag{A63.2}
\]
Every $e_pR$ here is finite $p$-primary torsion, since $v\nmid p$.
Thus $T_v$ is torsion and $Q_v$ is algebraic over $L$, even
though the coefficient $a_p$ need not be an algebraic number.
If $Q_v\in U(L_v)$, the local section is its point section and
the vanishing of the $q$-adic component in A.62 says that
\[
 \frac{f(Q_v)}{f(B)}\in
 \ker\bigl(L_v^\times\longrightarrow
                  \varprojlim_rL_v^\times/L_v^{\times q^r}\bigr).
\]
This kernel is the finite group of prime-to-$q$ roots of unity.
Indeed valuation contributes $\Z_q$, and an open subgroup of
the unit group is a finite-rank $\Z_q$-lattice; the remaining
kernel is exactly its prime-to-$q$ torsion. Since the displayed
ratio is algebraic, being a root of unity in $L_v$ makes it a
root of unity as an algebraic number.

Consider the closed irreducible graph curve in
$E\times\mathbf G_m$,
\[
 \mathcal C_f=
 \left\{\left(Q-B,\frac{f(Q)}{f(B)}\right):
             Q\in E-\operatorname{supp}(\operatorname{div}f)\right\}.
 \tag{A63.3}
\]
It is closed because it is the graph of the morphism
$f:E\to\mathbf P^1$, intersected with
$E\times\mathbf G_m$, followed by translation and scaling.
Equation (A63.2) and the preceding root-of-unity conclusion make
its point associated with $Q_v$ a torsion point of the product.

The Manin--Mumford theorem for a product of an abelian variety
and a torus applies to precisely this group and this curve
\citep[Theorem~1, \S4, p.~587]{AMHindry1988}. Its conclusion is
that the torsion points of a closed subvariety lie in finitely
many torsion translates of algebraic subgroups contained in
that subvariety. The curve (A63.3) contains no positive-dimensional
such translate. Its projection to $E$ is nonconstant of degree
one. A one-dimensional connected subgroup with nonzero
projection to $E$ is an elliptic curve, since every homomorphism
$\mathbf G_m\to E$ is zero; its projection to $\mathbf G_m$
is then constant, since every morphism from a proper elliptic
curve to $\mathbf G_m$ is constant. This contradicts the
nonconstancy of $f$. Thus the set of torsion points in
$\mathcal C_f$ is finite.

Let $Z\subset E$ be the finite reduced $L$-subscheme consisting
of their inverse images and of $E[N]$. The torsion intersection
is Galois stable, so it does define such a subscheme.
At every place under consideration, an interior $Q_v$ belongs
to $Z$ by (A63.3), and a cusp centre belongs to $Z$ by
construction. Therefore the proper adelic point $aR$ is
supported on $Z$ at all but finitely many places.

This adelic point lies in the closure of $E(L)$: the map
$m\mapsto mR$ extends continuously from $\Z$ to
$\widehat\Z$ on the product of the local point groups at finite places
and the archimedean component groups. Hence it survives finite
abelian descent, whose locus is closed and contains $E(L)$.
Apply \citet[Theorem~8.2]{AMStoll2007}, with $C=E$ of genus one
and the finite subscheme $Z$, using only the density-one set
of finite places. It gives one point $z\in Z(L)$ equal to this
modified adelic point. This use is unaffected by the real-place
correction \citep{AMStollErrata2017}.

Finally this equality of local points determines the exponent.
Choose a coefficient prime $\ell\notin\mathcal P$ and a place
above it. The $\ell$-adic Kummer class of $z-B$ is zero there,
so $z-B$ is prime-to-$\ell$ torsion in that local point group.
It follows that $z-B$ is a global torsion point. Repeating over
all $\ell\notin\mathcal P$ shows that its order is supported
on $\mathcal P$. For each $p\in\mathcal P$, at a place above
$p$ take the local elliptic logarithm. It kills torsion and has
nonzero value on $R$: its kernel is torsion, since on a
sufficiently small formal subgroup it is injective. From the
local equality $z=aR$ we obtain
\[
 0=\log(z-B)=(a_p-n)\log R,
\]
and consequently $a_p=n$. This proves (A63.1).
The existence for integer exponents is supplied by the indicated
actual point and cusp sections, which already exist in the full
fundamental group.
\end{proof}

For A.12 the curve $E$ has $j=0$, so every elliptic twist in
that question satisfies the CM hypothesis after a finite
extension. The argument therefore proves the same mechanism
for every such twist, relative to a chosen flex origin and a
nontorsion point. It uses the full geometric abelian quotient
and genuine all-place Selmerness. The cuspidally central
abstract lifts of A.59 still exist, but cannot be Selmer for
the excluded exponents. Unlike the special erasure proof
A.60, the argument uses unbounded toric coefficient primes;
it does not assert that one fixed finite coefficient quotient
detects all the exponents in A.63.

Exponents which differ from every integer at infinitely many
coefficient primes are not decided by A.63. Nor does this
identify the exact section set of a particular punctured
congruence cover. It does, however, give an exact answer on
the entire class of finite coefficient alterations, including
one altered prime, already before the intrinsic
metabelian level. The proofs do not require a rank computation
or finiteness of a Tate--Shafarevich group over the splitting field.
